\chapter{記号}
\label{sec:notation}

ニューロンの型は、主な神経伝達物質の名前の最初の4文字で表す。\nt{GLUT}：グルタミン酸、\nt{GABA}：$\gamma$-アミノ酪酸、\nt{GLYC}：グリシン、\nt{ACET}：アセチルコリン、\nt{DOPA}：ドーパミン、\nt{SERO}：セロトニン、\nt{HIST}：ヒスタミン、\nt{NORA}：ノルアドレナリン、\nt{ADRE}：アドレナリン、\nt{ASPA}：アスパラギン酸（表\ref{tab:types}）。

文字の数は量の数より少ないので、同じ文字が章によって別のものを表すことがある。下の表では、意味ごとに1行を割り当てた。右の列は、その記号を導入した式である。

\begin{center}
\footnotesize
\setlength{\tabcolsep}{4pt}
\renewcommand{\arraystretch}{1.12}
\begin{longtable}{>{\raggedright}p{2.0cm}>{\raggedright}p{9.6cm}>{\raggedright\arraybackslash}p{2.4cm}}
\toprule
記号 & 意味 & 導入箇所 \\
\midrule
\endfirsthead
\toprule
記号 & 意味 & 導入箇所 \\
\midrule
\endhead
\bottomrule
\endfoot
$a$ & 線形伝達特性の傾き、$f(u)=au$ & \eqref{eq:feedback} \\
$a$ & \nt{ACET} のレベル：$0$ は読み出し、$1$ は書き込み & \eqref{eq:acetnet} \\
$a_i$, $a$ & 群の疲労：リズム発生器で。徐波睡眠のシーソーで & \eqref{eq:halfcentre}, \eqref{eq:updown} \\
$a_S$, $a_P$ & 心拍リズムのアクセルとブレーキに対する感度 & \eqref{eq:gasbrake} \\
$A(r)$, $A_0$ & 頻度 $r$ でのスパイクに対するシナプスの応答。休んだシナプスでの応答 & \eqref{eq:depression} \\
$A_1$ & 伝達物質1包に対する受け手の応答 & \eqref{eq:quantal} \\
$A$ & 樹状突起を含むニューロンの膜面積 & \eqref{eq:dendriteload} \\
$A_f$, $A_s$ & 次の運動まで保持される補正の割合（速い過程と遅い過程） & \eqref{eq:tworate} \\
$b_i$ & ペースメーカー自身の流入 & \eqref{eq:seroneuron} \\
$b$ & リズム発生器や睡眠のシーソーで、働きが群を疲れさせる速さ & \eqref{eq:halfcentre}, \eqref{eq:updown} \\
$b$ & デジタル化の1サンプルあたりのビット数 & \eqref{eq:probedata} \\
$B_f$, $B_s$ & 補正が誤差から学習する強さ（速い過程と遅い過程） & \eqref{eq:tworate} \\
$c$ & 空間内の伝達物質の濃度 & \eqref{eq:seroneuron} \\
$c$ & 隣り合うニューロンのノイズの相関係数 & \eqref{eq:correlated} \\
$c$ & カウントの差がラケットを動かす係数 & \eqref{eq:dishreadout} \\
$c_f$ & $C$ ニューロンの応答：すべての位置にわたる $s_{f,p}$ の最大値 & \eqref{eq:scstage} \\
$c_m$ & 膜の比容量 & \eqref{eq:dendriteload} \\
$C$ & 努力のコスト：時間 $\tau_a$ での行動のコストは $C/\tau_a$ & \eqref{eq:vigor} \\
$d'$ & 二つの入力の弁別度：差をノイズで割ったもの & \eqref{eq:gainsnr} \\
$d$, $d_j$ & 信号経路上の遅延 & \eqref{eq:glutneuron}, \eqref{eq:vstar} \\
$d$ & 筋肉とのフィードバックループの遅延 & \eqref{eq:delaystable} \\
$D$ & 拡散係数 & \eqref{eq:diffusionlength} \\
$D$ & \nt{GLUT}--\nt{GABA} ネットワークの頻度の式の分母 & \eqref{eq:isn} \\
$D$ & 遅延報酬の待ち時間 & \eqref{eq:patience} \\
$e$, $e_{ij}$ & シナプスの適格度トレース & \eqref{eq:threefactor} \\
$e$, $e_i$ & 誤差配線を通って届く、出力や運動の誤差 & \eqref{eq:lms}, \eqref{eq:adaptivefilter}, \eqref{eq:tworate} \\
$E_{\text{spk}}$ & シナプスの働きを含むスパイク1回のエネルギー & \eqref{eq:energy} \\
$E$ & \nt{GLUT} 群の頻度 & \eqref{eq:ei} \\
$E_E$ & 興奮性シナプスの平衡電位 & \eqref{eq:shunt} \\
$f$, $F$ & 伝達特性：入力の関数としての頻度 & \eqref{eq:ratenet}, \eqref{eq:gain} \\
$f$, $f_0$ & 心拍数。ペースメーカー自身のリズム & \eqref{eq:gasbrake} \\
$f_s$ & 電極1本のサンプリング周波数 & \eqref{eq:probedata} \\
$F(t)$, $F_1$ & 筋力。単収縮1回の力 & \eqref{eq:forcesum} \\
$F_1$, $F_2$ & 関節を逆向きに引く二つの筋肉の力 & \eqref{eq:antagonists} \\
$g$ & 受容体の濃度に対する感度 & \eqref{eq:seroneuron} \\
$g_L$, $g_E$, $g_I$ & 漏れ、興奮、抑制のコンダクタンス & \eqref{eq:shunt} \\
$g$ & \nt{NORA} が設定するゲイン & \eqref{eq:gain}, \eqref{eq:machine} \\
$g$ & 記憶ネットワークでの \nt{GABA} による全体の抑制 & \eqref{eq:acetnet} \\
$g$ & 上から設定される痛みのゲートのゲイン & \eqref{eq:gate} \\
$g$ & ホルモンカスケード1段のゲイン & \eqref{eq:cascade} \\
$g_j$ & 適応フィルタの入力側にある、固有の時定数をもつフィルタ $j$ & \eqref{eq:adaptivefilter} \\
$G$ & フィードバックループのゲイン & \eqref{eq:feedback} \\
$G$ & 遅延のあるループでの修正の速さ & \eqref{eq:delaystable} \\
$h$, $h_I$ & \nt{GLUT} 群と \nt{GABA} 群への外部入力 & \eqref{eq:ei}, \eqref{eq:isn} \\
$h(t)$ & 筋肉の単収縮1回 & \eqref{eq:twitch} \\
$H(t)$ & 睡眠圧の上側の閾値：ここに達するとマシンは眠る & \eqref{eq:sleeppressure} \\
$I$, $I_i$ & ニューロンや群への外部からの流入 & \eqref{eq:ratenet}, \eqref{eq:wta}, \eqref{eq:updown} \\
$I$ & \nt{GABA} 群の頻度 & \eqref{eq:ei} \\
$I$ & 刺激の明るさや強度 & \eqref{eq:photonnoise}, \eqref{eq:nakarushton} \\
$I$ & 膜を充電する入力電流 & \eqref{eq:dendriteload} \\
$k$ & バースト内のスパイク数 & \eqref{eq:burst} \\
$k$ & 増分がノイズの何倍でなければならないか & \eqref{eq:photonnoise} \\
$k$ & ホルモンカスケードのフィードバック係数 & \eqref{eq:cascade} \\
$k_s$ & 痛みが感作を押し上げる強さ & \eqref{eq:sensitization} \\
$K$ & 結果の特徴の数 & \eqref{eq:vectortd} \\
$K$ & 関節の剛性 & \eqref{eq:antagonists} \\
$K$ & ホルモンカスケードのループゲイン、$K=g^3k$ & \eqref{eq:cascade}, \eqref{eq:cascadestable} \\
$K_\varphi$ & 概日発振器が外部信号に同調する強さ & \eqref{eq:pll} \\
$L$ & 遅延線の長さ & \eqref{eq:itd} \\
$L$ & 痛みのセンサーからの配線の長さ & \eqref{eq:paintime} \\
$L(t)$ & 睡眠圧の下側の閾値：ここに達するとマシンは目覚める & \eqref{eq:sleeppressure} \\
$M$ & 予兆と報酬の因果関係の尺度 & \eqref{eq:retro} \\
$M_k$ & 特徴 $k$ の期待値 & \eqref{eq:vectortd} \\
$M_{ik}$ & \nt{GABA} ニューロン $k$ からの抑制の強さ & \eqref{eq:machine} \\
$M$ & 関節のトルク & \eqref{eq:antagonists} \\
$M$ & ミキサーの入力線の数 & \eqref{eq:cover} \\
$n$, $\bar n$ & 窓内のスパイク数。その平均 & \eqref{eq:rateerror} \\
$n$ & 閾値素子の入力数 & \eqref{eq:cover} \\
$n_\uparrow$, $n_\downarrow$ & 「上」と「下」の領域の窓あたりのスパイク数 & \eqref{eq:dishreadout} \\
$N$ & ニューロン数 & \eqref{eq:correlated}, \eqref{eq:energy} \\
$N$ & プローブの電極数 & \eqref{eq:probedata} \\
$N_s$ & 2個のニューロンの間の放出部位の数 & \eqref{eq:quantal} \\
$p$ & スパイク1回あたりの伝達物質放出確率 & \eqref{eq:burst} \\
$p$ & 記憶ネットワークで活動しているニューロンの割合 & \eqref{eq:acetnet} \\
$p$ & 補償を学習する外乱 & \eqref{eq:tworate} \\
$P$ & 信号に使われる電力 & \eqref{eq:energy}, \eqref{eq:dishpower} \\
$P$ & 保持している推定値の不確かさ & \eqref{eq:kalman} \\
$P$ & 「ブレーキ」の活動：心臓への \nt{ACET} 配線 & \eqref{eq:gasbrake} \\
$P_{\max}$ & 閾値素子が二つのクラスに分けられるランダムなパターンの数 & \eqref{eq:cover} \\
$P_{\text{miss}}$ & ミスの確率 & \eqref{eq:dishdensity} \\
$q$ & 世界が変わる速さ & \eqref{eq:kalman} \\
$q_k$ & \nt{GABA} ニューロン $k$ の頻度 & \eqref{eq:machine} \\
$q_0$ & 抑制性介在ニューロンの背景活動 & \eqref{eq:gate} \\
$q$ & 痛みのゲートにおける抑制性介在ニューロンの頻度 & \eqref{eq:gate} \\
$q$ & 次の運動単位が前の運動単位の何倍強いか & \eqref{eq:sizeprinciple} \\
$r$, $r_i$ & ニューロンのスパイク頻度 & \eqref{eq:ratenet} \\
$r$, $r_t$ & 報酬（報酬の流れ） & \eqref{eq:rpe}, \eqref{eq:ctd} \\
$\mathbf r(t)$ & リザバーの応答ベクトル & \eqref{eq:reservoir} \\
$R$, $\bar R$ & 得られた報酬。その期待値、または単位時間あたりの平均 & \eqref{eq:perturbation}, \eqref{eq:vigor} \\
$R$ & フィルタ入力のノイズの分散 & \eqref{eq:kalman} \\
$R$, $r_0$ & 遅延報酬と即時報酬 & \eqref{eq:patience} \\
$R_{\max}$ & 最大伝送速度、bit/s & \eqref{eq:spikeentropy} \\
$r_{\max}$ & 最大スパイク頻度 & \eqref{eq:gain}, \eqref{eq:nakarushton} \\
$R_{\text{raw}}$, $R_{\text{spike}}$ & プローブのデータ量：生データとスパイクのみ、bit/s & \eqref{eq:probedata} \\
$s_j$ & ニューロン $j$ の出力の符号 & \eqref{eq:dalesign} \\
$s_i$ & 受け手の受容体の符号 & \eqref{eq:seroneuron} \\
$s$, $\hat s$ & 世界の状態。その推定値 & \eqref{eq:rpe}, \eqref{eq:kalman} \\
$s_{f,p}$ & 位置 $p$ における特徴 $f$ への $S$ ニューロンの応答 & \eqref{eq:scstage} \\
$S$ & 「アクセル」の活動：心臓への \nt{NORA} 配線 & \eqref{eq:gasbrake} \\
$S$ & 睡眠圧 & \eqref{eq:sleeppressure} \\
$t_1$ & 最初のスパイクの時刻 & \eqref{eq:latency} \\
$t_k$ & $k$ 番目のスパイクの時刻 & \eqref{eq:forcesum} \\
$t_{f,i}$ & 特徴 $f$ のテンプレート & \eqref{eq:scstage} \\
$T$ & スパイクを数える窓。光子の蓄積時間 & \eqref{eq:rateerror}, \eqref{eq:photonnoise} \\
$T_c$ & 筋肉の単収縮の立ち上がり時間 & \eqref{eq:twitch} \\
$T_{\text{burst}}$ & 培養の一斉発火の間隔 & \eqref{eq:burstinterval} \\
$u$ & ニューロンへの入力の総和 & \eqref{eq:gain} \\
$u$, $u(t)$ & 脳の指令：適応フィルタの入力として。ホルモンカスケードへの指令として & \eqref{eq:adaptivefilter}, \eqref{eq:cascade} \\
$u(t)$ & リザバーの入力 & \eqref{eq:reservoir} \\
$U$ & 一定の入力のレベル & \eqref{eq:latency} \\
$U$ & スパイク1回で放出される伝達物質の蓄えの割合 & \eqref{eq:depression} \\
$v$ & 配線中の信号の速さ。光点の動く速さ & \eqref{eq:itd}, \eqref{eq:vstar} \\
$V$ & 膜電圧 & \eqref{eq:glutneuron}, \eqref{eq:shunt} \\
$V$ & 将来の報酬の期待値 & \eqref{eq:rpe}, \eqref{eq:ctd} \\
$w$, $w_{ij}$, $W_{ij}$ & 結合の重み & \eqref{eq:glutneuron}, \eqref{eq:ratenet}, \eqref{eq:acetnet}, \eqref{eq:lms}, \eqref{eq:adaptivefilter} \\
$\mathbf w_k$ & $C$ ニューロン $k$ の入力の重み & \eqref{eq:traceinvariance} \\
$w_{q\mu}$, $w_{q\nu}$, $w_{z\mu}$ & 痛みのゲートの結合の重み & \eqref{eq:gate} \\
$W_{\text{out}}$ & リザバーの線形読み出しの重み & \eqref{eq:reservoir} \\
$x$, $y$ & 論理入力と論理出力 & \eqref{eq:dualrail}, \eqref{eq:veto} \\
$x$, $y$ & 送り手と受け手の活動 & \eqref{eq:hebb} \\
$x$ & 遅延線上の検出器の位置 & \eqref{eq:itd} \\
$x_i$ & 感覚器からの入力 & \eqref{eq:acetnet} \\
$x_p$ & 位置 $p$ の入力 & \eqref{eq:scstage} \\
$\mathbf x$ & $S$ ニューロンの出力、$C$ ニューロンの入力 & \eqref{eq:traceinvariance} \\
$x_j$ & 適応フィルタの入力 $j$：フィルタ $g_j$ を通った指令 & \eqref{eq:lms}, \eqref{eq:adaptivefilter} \\
$x_f$, $x_s$ & 速い過程と遅い過程の補正 & \eqref{eq:tworate} \\
$x_1$, $x_2$, $x_3$ & ホルモンカスケード3段のレベル。$x_3$ は最終ホルモン & \eqref{eq:cascade} \\
$x^*$ & 新しい雪崩が始まる、伝達物質の蓄えの回復割合 & \eqref{eq:burstinterval} \\
$y$ & フィルタのノイズを含む入力 & \eqref{eq:kalman} \\
$y$, $\hat y$ & センサーの信号。適応フィルタによるその予測 & \eqref{eq:adaptivefilter} \\
$y_k$, $\bar y_k$ & $C$ ニューロン $k$ の活動。そのトレース & \eqref{eq:traceinvariance} \\
$y(t)$ & リザバーの読み出しの出力 & \eqref{eq:reservoir} \\
$z$ & 痛みのゲートの出力の頻度 & \eqref{eq:gate} \\
\midrule
$\alpha_i^\pm$ & 正の誤差と負の誤差の重み & \eqref{eq:asymmetric} \\
$\beta$ & 相互抑制の強さ & \eqref{eq:wta} \\
$\beta$ & 痛みのゲートの出力に対する抑制の強さ & \eqref{eq:gate} \\
$\gamma$ & 割引率 & \eqref{eq:rpe} \\
$\delta(\cdot)$ & デルタ関数：瞬間的な跳び & \eqref{eq:glutneuron} \\
$\delta$, $\delta_t$ & 報酬予測誤差 & \eqref{eq:rpe}, \eqref{eq:ctd} \\
$\delta t$ & 両耳への音の到達時間差 & \eqref{eq:itd} \\
$\Delta$, $\Delta_{\max}$ & スパイクの間隔。同時性の窓 & \eqref{eq:window} \\
$\Delta t$ & 受け手が時間を区別できる精度 & \eqref{eq:spikeentropy} \\
$\Delta F$ & 次の運動単位による力の増分 & \eqref{eq:sizeprinciple} \\
$\Delta I$ & 区別できる明るさの増分 & \eqref{eq:photonnoise} \\
$\Delta p$ & 窓あたりのラケットの移動 & \eqref{eq:dishreadout} \\
$\Delta u$ & 二つの群の入力の差 & \eqref{eq:gainsnr} \\
$\Delta V$ & 膜電圧をどれだけ上げる必要があるか & \eqref{eq:dendriteload} \\
$\Delta x$ & 隣り合うセンサーの間隔 & \eqref{eq:vstar} \\
$\varepsilon$ & 入力の相対的な差 & \eqref{eq:wtagain} \\
$\eta$ & 学習率 & \eqref{eq:plasticity}, \eqref{eq:lms}, \eqref{eq:traceinvariance} \\
$\eta$ & 増幅器の後のネットワーク内のノイズ & \eqref{eq:softmax} \\
$\mu$ & 痛みのゲートへの触覚入力 & \eqref{eq:gate} \\
$\nu$ & 痛みの入力 & \eqref{eq:gate} \\
$\theta$ & 発火の閾値 & \eqref{eq:window}, \eqref{eq:scstage} \\
$\kappa$ & スパイク1回あたりの伝達物質の量 & \eqref{eq:seroneuron} \\
$\kappa_i$ & 正の誤差と負の誤差の重みの比 & \eqref{eq:expectile} \\
$\kappa_m$ & 筋力と筋肉の剛性の関係 & \eqref{eq:antagonists} \\
$\ell$ & 伝達物質が広がる距離 & \eqref{eq:diffusionlength} \\
$\ell$ & 筋肉が関節を引くモーメントアーム & \eqref{eq:antagonists} \\
$\xi$ & ニューロン出力のランダムな揺らぎ & \eqref{eq:perturbation} \\
$\rho(\boldsymbol\psi)$ & 設定 $\boldsymbol\psi$ で過ごす時間の割合 & \eqref{eq:dishdensity} \\
$\sigma$ & ばらつき、ノイズの尺度 & \eqref{eq:rateerror}, \eqref{eq:softmax} \\
$\sigma$ & 応答が最大の半分になる明るさ & \eqref{eq:nakarushton} \\
$\sigma_{\text{pre}}$, $\sigma_{\text{post}}$ & 増幅器の前と後のノイズ & \eqref{eq:gainsnr} \\
$\tau$ & 膜の時定数 & \eqref{eq:glutneuron} \\
$\tau$ & ホルモンカスケード1段の時定数 & \eqref{eq:cascade} \\
$\tau_{\text{eff}}$ & 自己興奮する群の時定数 & \eqref{eq:perfectint} \\
$\tau_c$ & 空間内の伝達物質の寿命 & \eqref{eq:seroneuron} \\
$\tau_E$, $\tau_I$ & \nt{GLUT} 群と \nt{GABA} 群の時定数 & \eqref{eq:ei} \\
$\tau_e$ & 適格度トレースの寿命 & \eqref{eq:threefactor} \\
$\tau_{\text{tr}}$ & $C$ ニューロンの活動のトレースの寿命 & \eqref{eq:traceinvariance} \\
$\tau_s$ & 損傷後に感度が戻るまでの時間 & \eqref{eq:sensitization} \\
$\tau_r$, $\tau_d$ & 覚醒中に睡眠圧がたまる時間と、睡眠中に解消される時間 & \eqref{eq:sleeppressure} \\
$\tau_{\text{rec}}$ & 伝達物質の蓄えの回復時間 & \eqref{eq:depression} \\
$\tau_\gamma$ & 計画の視野 & \eqref{eq:ctd} \\
$\tau_v$ & 期待値が更新される速さ & \eqref{eq:ramp} \\
$\tau_a$ & リズム発生器や睡眠のシーソーにおける疲労の時間 & \eqref{eq:halfcentre}, \eqref{eq:updown} \\
$\tau_a^*$ & 行動を実行する最適な時間 & \eqref{eq:vigor} \\
$\Phi$ & 学習則の右辺 & \eqref{eq:plasticity} \\
$\Phi$ & 光子の流束 & \eqref{eq:photonnoise} \\
$\varphi_k$ & 得られた結果の特徴 $k$ & \eqref{eq:vectortd} \\
$\varphi$ & 概日発振器と外部信号の位相差 & \eqref{eq:pll} \\
$\boldsymbol\psi$ & DishBrain モデルにおけるネットワークの設定 & \eqref{eq:dishdensity} \\
$\omega$, $\omega_0$ & 外部信号（光）の周波数。概日発振器の固有周波数 & \eqref{eq:pll} \\
\end{longtable}
\end{center}
